\subsection{GREATEST ELEMENT AND LEAST ELEMENT}

If there exists an element \(a\) in an ordered set \(\mathbf{E}\) such that \(a \leqslant x\) for all \(x \in \mathbf{E}\) , then \(a\) is the only element of \(\mathbf{E}\) with this property; for if also \(b \leqslant x\) for all \(x \in \mathbf{E}\) , then we should have \(a \leqslant b\) and \(b \leqslant a\) , and consequently \(a = b\) .

DEFINITION 4. Let E be an ordered set. An element \(a \in E\) is said to be the least (resp. greatest) element of E if for all \(x \in E\) we have \(a \leqslant x\) (resp. \(x \leqslant a\) ).

An ordered set need not have a greatest element nor a least element. If E has a least element a, then a is the greatest element with respect to the opposite ordering.

If E has a least element \(a\) , then \(a\) is the unique minimal element of E; for if \(x \in E\) is distinct from \(a\) , we have \(a < x\) .

Examples

\begin{enumerate}
\def\labelenumi{(\arabic{enumi})}
\item
  Let \(\mathfrak{S}\) be a non-empty subset of the set \(\mathfrak{P}(\mathbf{E})\) of subsets of a set \(\mathbf{E}\) . If \(\mathfrak{S}\) has a least (resp. greatest) element A with respect to the inclusion relation, then A is the intersection (resp. union) of the sets of \(\mathfrak{S}\) . Conversely, if the intersection (resp. union) of the sets of \(\mathfrak{S}\) belongs to \(\mathfrak{S}\) , then it is the least (resp. greatest) element of \(\mathfrak{S}\) .
\item
  In particular, \(\varnothing\) is the least element and E the greatest element of \(\mathfrak{P}(E)\) . In the set \(\Phi(E, F)\) of mappings of subsets of E into F, ordered by extension of mappings (no. 1, Example 3), the empty mapping is the least element, and there is no greatest element unless F consists of a single element. The diagonal \(\Delta\) of \(E \times E\) is the least element of the set of graphs of equivalence relations on E (or of the set of preorderings on E).
\end{enumerate}

PROPOSITION 3. Let E be an ordered set and let E' be the disjoint union of E and a set \{a\} consisting of a single element. Then there exists a unique ordering on E' which induces the given ordering on E and for which a is the greatest element of E'.

For if G is the graph of the ordering on E, the graph of an ordering on \(E'\) which satisfies the conditions of the Proposition must be the union \(G'\) of G and the set of all pairs \((x, a)\) where \(x \in E'\) ; conversely, it is clear that \(G'\) is the graph of an ordering on \(E'\) which satisfies the given conditions. ¶ The ordered set \(E'\) is said to be obtained by adjoining a greatest element a to E (cf.~Exercise 3).

A subset A of a preordered set E is said to be cofinal (resp. coinitial) in E if for every \(x \in E\) there exists \(y \in A\) such that \(x \leqslant y\) (resp. \(y \leqslant x\) ). To say that an ordered set E has a greatest (resp. least) element therefore means that E has a cofinal (resp. coinitial) subset consisting of a single element.

